\documentclass[conference]{IEEEtran}
\usepackage{amsmath,amssymb}
\title{Uncertainty-Aware Scheduling for Clinical Decision Support}
\author{Author A}
\begin{document}
\maketitle

\begin{abstract}
Clinical decision support systems must schedule inference under uncertain patient arrival and uncertain model calibration.
We formulate the scheduling problem with a calibration-aware objective and propose a policy that reserves capacity for high-uncertainty cases.
On a retrospective dataset, the policy improves the fraction of high-uncertainty cases served within the clinical deadline by 19\%.
\end{abstract}

\section{Introduction}
\label{sec:intro}
Hospitals increasingly rely on models that flag deterioration risk.
These models are calibrated on historical cohorts and drift when case mix changes.
A scheduler that ignores calibration drift may allocate capacity to cases the model already handles well.
We therefore treat calibration uncertainty as a scheduling signal.
Our contributions are a calibration-aware objective, a reservation policy, and a retrospective evaluation.

\section{Related Work}
\label{sec:related}
Prior work on clinical scheduling optimizes waiting time and utilization.
Prior work on calibration studies reliability diagrams and recalibration methods.
The two literatures rarely meet, although both concern uncertainty.
Our formulation connects them through a per-case uncertainty score.

\section{Problem Formulation}
\label{sec:problem}
Let $u_i$ denote the uncertainty score of case $i$ and let $d_i$ denote its deadline.
The scheduler chooses a service order to maximize the number of high-uncertainty cases completed before $d_i$, subject to capacity $C$.
We write the objective as
\begin{equation}
\label{eq:clinical}
\max \sum_i \mathbf{1}[u_i \ge u_0] \, \mathbf{1}[T_i \le d_i].
\end{equation}

\section{Method}
\label{sec:method}
We reserve a fraction $\beta$ of capacity for cases with $u_i \ge u_0$ and serve the remainder in deadline order.
The reservation fraction is adjusted weekly from observed calibration error.
In particular, the adjustment uses only aggregate error, so it does not require patient-level labels at deployment.

\section{Evaluation}
\label{sec:eval}
We replay a retrospective cohort of 12{,}000 cases.
The reservation policy improves on-time service for high-uncertainty cases by 19\%.
Overall utilization drops by 3\%, which the clinical team considered acceptable.
A sensitivity study shows that the gain persists when the uncertainty threshold varies by 20\%.

\section{Conclusion}
\label{sec:conclusion}
Calibration uncertainty is a useful scheduling signal.
A simple reservation policy captures most of the benefit.
Future work will evaluate the policy prospectively and study fairness across demographic groups.

\end{document}
